\documentclass[11pt]{article}

% ---------- Packages ----------
\usepackage[margin=1in]{geometry}
\usepackage{amsmath,amssymb,amsthm}
\usepackage{enumitem}
\usepackage{booktabs}
\usepackage{longtable}
\usepackage{array}
\usepackage{xcolor}
\usepackage{titlesec}
\usepackage{fancyhdr}
\usepackage[colorlinks=true,linkcolor=blue!50!black,citecolor=blue!50!black,urlcolor=blue!50!black]{hyperref}
\usepackage{parskip}
\setlength{\emergencystretch}{3em}

% ---------- Colors ----------
\definecolor{hubblue}{RGB}{31,73,125}
\definecolor{extgray}{RGB}{110,110,110}

% ---------- Section styling ----------
\titleformat{\section}{\normalfont\Large\bfseries\color{hubblue}}{\thesection}{1em}{}
\titleformat{\subsection}{\normalfont\large\bfseries\color{hubblue}}{\thesubsection}{1em}{}
\titleformat{\subsubsection}{\normalfont\bfseries}{\thesubsubsection}{1em}{}

\pagestyle{fancy}
\fancyhf{}
\fancyhead[L]{\small\itshape Chapter-1 Computational Companion}
\fancyhead[R]{\small\itshape\rightmark}
\fancyfoot[C]{\thepage}
\renewcommand{\headrulewidth}{0.3pt}

% ---------- Custom commands ----------
\newcommand{\extnote}[1]{{\small\color{extgray}\emph{[External / not in Pavliotis: #1]}}}
\newcommand{\pav}[1]{{\small\color{hubblue}(Pavliotis, #1)}}
\newcommand{\Dstar}{D_{\star}}

% Feature-brief field list (the master-plan altitude replacement for the old
% per-unit implementation-contract tables, which were deleted: call signatures,
% code, and tolerances are owned downstream, not here).
\newenvironment{brief}{%
  \begin{description}[style=sameline,leftmargin=2.3cm,font=\bfseries\color{hubblue}]
}{%
  \end{description}
}

\title{\textbf{A Computational Companion to Pavliotis, Chapter 1}\\[4pt]
\large Master Plan: Covariance Operators, Spectral Theory, and the Karhunen--Lo\`eve Expansion}
\author{Stefan}
\date{\today}

\begin{document}

\maketitle

\begin{abstract}
\noindent This is the \emph{master plan} for a computational project in Julia that deepens
understanding of Chapter~1 of Pavliotis, \emph{Stochastic Processes and Applications}, by
implementing its core objects and verifying each against an analytic prediction. It owns the
project's through-line, its decomposition into feature briefs, the repository architecture and its
limits, the cross-cutting conventions, a consolidated risk register, and the budget. It
deliberately does \emph{not} carry call signatures, code, or numerical tolerances: those are
resolved downstream --- in the per-commit plans that each feature brief is handed to, and in the
shipped code itself --- and a copy here would be a second, drifting source of truth. The project is
organized around a single hub, the covariance operator $\mathcal{C}$, with two diagonalizing bases
(Fourier/Bochner and Mercer/Karhunen--Lo\`eve) and \emph{four} mathematical square roots used for
sampling (Cholesky, KL truncation, random Fourier series, circulant embedding) --- of which
\emph{three are implemented} (Cholesky, KL, circulant); the random-Fourier series is described as
structure but not built, and this costs no capability because circulant embedding is itself the
Bochner square root. Ergodicity closes the loop between one simulated path and the analytic objects
it is checked against. The plan excludes machine-learning and SciML content. At the working budget
of \textbf{15\,h/week} ($6\times 2.5$\,h) the full planned ladder is \textbf{Units 0--7}: the
covariance/spectral spine (Units 0--4), Brownian motion as a scaling limit (Unit~5), fractional
Brownian motion (Unit~6), and a forward-pointing SDE bridge to Chapter~3 (Unit~7). \textbf{Units
0--3 are shipped}; Units~4--5 are committed core; Unit~6 then Unit~7 are the cut order if the buffer
is consumed. Each unit is written as a feature brief --- its one idea, the claim it proves and the
exact analytic target it is checked against, the falsifier it is designed to exhibit, its place in
the build DAG, the repository areas it touches, its effort and dominant risk, and its status.
Statements drawn from the text are cited by section/theorem number; statements standard but external
to Pavliotis are flagged explicitly. All citations are to the project's copy of the text, the
\emph{November 11, 2015 draft} (numbering differs from the published edition), checked directly
against it.
\end{abstract}

\tableofcontents

% =====================================================================
\section{The Through-Line}
% =====================================================================

\subsection{The hub: the covariance operator}

Every object studied in Chapter~1 is a fact about a single operator. For a mean-zero, second-order
process $X_t$, $t\in[0,T]$, with continuous covariance $R(t,s) = \mathbb{E}(X_tX_s)$, define the
integral operator
\begin{equation}
(\mathcal{C}f)(t) \;:=\; \int_0^T R(t,s)\,f(s)\,ds, \qquad f \in L^2[0,T].
\label{eq:covop}
\end{equation}
\pav{eqn.~1.33} This is self-adjoint, nonnegative, and (for continuous $R$) compact and
Hilbert--Schmidt \pav{Exercises 21--22}. Everything downstream is one of two operations performed on
$\mathcal{C}$.

\subsection{Two diagonalizations}

\paragraph{Fourier / Bochner (gated by stationarity).}
If $X_t$ is weakly stationary, $R(t,s) = R(t-s)$ is a function of a single lag, and Bochner's theorem
applies:
\begin{equation}
R(\tau) \;=\; \int_{\mathbb{R}} e^{i\omega\tau}\,\rho(d\omega),
\label{eq:bochner}
\end{equation}
for a unique nonnegative measure $\rho$ on $\mathbb{R}$ with $\rho(\mathbb{R}) = R(0)$
\pav{Thm.~1.14}. When $\rho$ is absolutely continuous, $\rho(d\omega) = S(\omega)\,d\omega$ with
\begin{equation}
S(\omega) = \frac{1}{2\pi}\int_{-\infty}^{\infty} e^{-i\omega\tau}R(\tau)\,d\tau,
\qquad
R(\tau) = \int_{-\infty}^{\infty} e^{i\omega\tau}S(\omega)\,d\omega.
\label{eq:spectraldensity}
\end{equation}
\pav{eqns.~1.7--1.8} $S(\omega)$ is the \emph{spectral density}; it lives on the frequency axis, not
on the sample space --- a distinction worth stating explicitly since both objects are conventionally
denoted $\omega$. Equations~\eqref{eq:spectraldensity} form a Fourier transform pair and invert by
the Fourier inversion theorem, independently of any symmetry of $R$. Realness and evenness of $R$ add
only that $S$ is itself real and even --- so the sign of $i$ in the exponent is immaterial (one may
replace $e^{\pm i\omega\tau}$ by $\cos\omega\tau$); the $1/2\pi$ is placed on the $S$-side purely by
convention. With this convention $\int_{\mathbb{R}} S(\omega)\,d\omega = R(0)$ (total variance),
which fixes the normalization every estimator in Unit~1 must reproduce.

\paragraph{Mercer / Karhunen--Lo\`eve (general, needs a compact domain).}
Independently of stationarity, the spectral theorem for compact self-adjoint operators gives a
countable orthonormal eigenbasis $\{e_k\}$ of $\mathcal{C}$ with eigenvalues $\lambda_k \geq 0$,
\begin{equation}
\mathcal{C}e_k = \lambda_k e_k, \qquad R(t,s) = \sum_{k=1}^{\infty}\lambda_k\,e_k(t)\,e_k(s),
\label{eq:kl-eigen}
\end{equation}
\pav{eqn.~1.32; Mercer's theorem, stated inline (unnumbered) in \S1.5} and the Karhunen--Lo\`eve
expansion
\begin{equation}
X_t = \sum_{k=1}^{\infty} \xi_k\, e_k(t), \qquad \xi_k = \int_0^T X_t\, e_k(t)\,dt,\quad
\mathbb{E}(\xi_k\xi_\ell) = \lambda_k\delta_{k\ell},
\label{eq:kl}
\end{equation}
\pav{Thm.~1.28 (Karhunen--Lo\`eve)} converging in $L^2$, uniformly in $t$. If $X_t$ is Gaussian, the
$\xi_k$ are independent (uncorrelated \emph{and} jointly Gaussian), and
\begin{equation}
X_t = \sum_{k=1}^{\infty} \sqrt{\lambda_k}\,\xi_k\, e_k(t), \qquad \xi_k \overset{\text{iid}}{\sim} N(0,1).
\label{eq:kl-gaussian}
\end{equation}
\pav{eqn.~1.36} The eigenvalues are simultaneously (i) eigenvalues of the covariance operator and
(ii) variances of the independent stochastic coefficients in the expansion --- one fact, two
readings.

\paragraph{When the two diagonalizations coincide.}
The Fourier basis diagonalizes $\mathcal{C}$ exactly when $\mathcal{C}$ is a convolution operator,
which requires the domain to be a group under addition and the kernel to be translation-invariant.
The interval $[0,T]$ is not a group; the circle (torus) is. On the torus, $\mathcal{C}$ is a genuine
convolution, its eigenfunctions are forced to be the characters $e^{ikt}$, and
\begin{equation}
\lambda_k = \hat R(k) \qquad \text{(exactly)},
\label{eq:torus-coincidence}
\end{equation}
the $k$-th eigenvalue equals the $k$-th Fourier coefficient of $R$ --- the KL sequence and the
(discrete, atomic) spectral measure are the same object. On $[0,T]$ the boundary breaks translation
invariance: the KL eigenfunctions of Brownian motion, for instance, are Sturm--Liouville sine modes
with eigenvalues $\lambda_n \propto (2n-1)^{-2}$ \pav{Example 1.29}, not Fourier modes, and
\eqref{eq:torus-coincidence} fails. \extnote{the general asymptotic recovery $\lambda_k \to \hat
R(\omega_k)$ as $T\to\infty$ --- where $\hat R(\omega)=2\pi S(\omega)$ is the \emph{un-normalized}
operator symbol, since the eigenproblem $\int R(t,s)e(s)\,ds=\lambda e$ carries no $1/2\pi$ --- is a
Toeplitz-operator fact (Grenander--Szeg\H{o}), not covered in Pavliotis. Stated precisely, Szeg\H{o}
controls the eigenvalue \emph{distribution} --- linear statistics $\frac1n\sum_k f(\lambda_k)$
converge to the average of $f\circ \hat R$ over the frequency band --- and \emph{not} each eigenvalue
uniformly in $k$; the extreme eigenvalues, at the edges of the spectrum, are exactly where that
control is weakest. This asymptotic --- not an exact $[0,T]$ identity, and not a uniform-in-$k$ one
--- is precisely the object of the Unit-3 cross-check, and the distributional/uniform distinction is
why that cross-check is restricted to the resolved bulk of the spectrum rather than gated on a
$\max_k$ over all of it.}

\subsection{One factorization, four routes (three implemented)}

Sampling a Gaussian process is applying a square root of $\mathcal{C}$ to white noise. Four square
roots give four sampling algorithms, all reproducing the same \emph{law}; \textbf{three are
implemented} (Cholesky, KL truncation, circulant embedding). The random-Fourier series is described
here as genuine mathematical structure but is \emph{not built} --- and, as noted below, this leaves
no capability gap.
\begin{itemize}[leftmargin=1.4cm]
  \item \textbf{Cholesky (general; implemented):} $\Sigma = LL^\top$ on a grid, $X = Lz$, $z\sim
  N(0,I)$.
  \item \textbf{KL truncation (general; implemented):} \eqref{eq:kl-gaussian} truncated at $K$ terms.
  \item \textbf{Random Fourier series, Gaussian amplitudes (stationary only; \emph{described, not
  implemented}):} on a \emph{one-sided} frequency grid $\omega_k>0$,
  $X(t) = \sum_k \big[a_k\cos(\omega_k t) + b_k\sin(\omega_k t)\big]$ with $a_k,b_k
  \overset{\text{iid}}{\sim} N\!\big(0,\,2\,S(\omega_k)\,\Delta\omega\big)$. The factor $2$ folds the
  two-sided spectrum onto $\omega>0$: $\mathrm{Cov}(\tau)=\sum_k 2S(\omega_k)\Delta\omega\cos(\omega_k
  \tau)\to R(\tau)$. \emph{Gaussian} amplitudes make this a genuine square root of $\mathcal{C}$
  applied to white noise, so the finite-$K$ sample is exactly Gaussian; a random-\emph{phase} variant
  (fixed amplitude $\sqrt{4S(\omega_k)\Delta\omega}$, uniform phase) reproduces the covariance but
  \emph{not} the Gaussian law at finite $K$ (Gaussian only in the limit, by the CLT), and is
  therefore not a square root. (This last is a decision record --- the reason the random-phase
  variant is rejected as a sampler --- and is retained even though the whole route is unimplemented.)
  \item \textbf{Circulant embedding (stationary only; implemented):} FFT-diagonalization of the
  circularly-extended Toeplitz covariance (Wood--Chan / Dietrich--Newsam); the circulant eigenvalues
  are the DFT of the covariance sequence --- i.e.\ the \emph{discrete} spectral density --- so this
  is precisely the Bochner square root made computational, on the grid rather than on the line. Exact
  and $O(m\log m)$ when the embedding is nonnegative definite.
\end{itemize}
Cholesky and KL are general; the random-Fourier and circulant routes require stationarity. Each is a
different basis in which to take the square root of the same operator: Cholesky is the triangular
factor, KL is the eigenbasis, and random-Fourier / circulant embedding are the Bochner square roots
on the continuous frequency axis and the discrete FFT grid respectively. \textbf{Because those two
are the \emph{same} Bochner square root in two realizations, shipping only circulant loses no
capability} --- the Bochner square root is present; the unbuilt route is the continuous-axis twin of
a route that is built.

\subsection{Stationarity as gate, ergodicity as loop-closer}

Stationarity is a \emph{precondition} for the Fourier vertex: Brownian motion is not stationary (its
covariance $\min(t,s)$ is not a function of $t-s$), so it has no spectral density; it possesses only
the Mercer/KL vertex. Ergodicity is the mechanism that licenses every theoretical check in this
project: for a stationary process with $C\in L^1(0,\infty)$,
\begin{equation}
\lim_{T\to\infty}\mathbb{E}\left|\frac{1}{T}\int_0^T X_s\,ds - \mu\right|^2 = 0,
\label{eq:ergodic}
\end{equation}
\pav{Prop.~1.16, eqn.~1.12} with finite-$T$ rate (for $C\geq 0$)
\begin{equation}
\mathbb{E}\left|\frac{1}{T}\int_0^T (X_s-\mu)\,ds\right|^2 = \frac{2}{T^2}\int_0^T (T-u)\,C(u)\,du
\;\leq\; \frac{2}{T}\int_0^\infty C(u)\,du \;=\; \frac{2\Dstar}{T},
\label{eq:ergodic-rate}
\end{equation}
\pav{eqn.~1.13, Lemma 1.17} where the first equality is general (Lemma~1.17) but the inequality uses
$C\geq0$ (for sign-changing $C$ the honest bound is $\tfrac{2}{T}\int_0^\infty|C|$, and
$\Dstar=\int_0^\infty C$ coincides with it only in the sign-definite case; the asymptotic form
$\mathrm{Var}\sim 2\Dstar/T$ holds regardless of sign given integrability). Here we introduce the
\textbf{Green--Kubo (transport) coefficient}
\begin{equation}
\Dstar \;:=\; \int_0^\infty C(t)\,dt,
\qquad
\mathbb{E}\Big(\int_0^t X_s\,ds\Big)^2 \approx 2\Dstar t \ \ (t \gg 1).
\label{eq:greenkubo}
\end{equation}
\pav{eqns.~1.14--1.15} \textbf{Notation caution (carry $\alpha$).} Pavliotis denotes both the
transport coefficient $\int_0^\infty C$ (eqn.~1.14) \emph{and} the noise strength in the OU kernel via
$C(0)=D/\alpha$ (Example~1.15) by the single symbol $D$. Example~1.18 makes the collision explicit:
there Pavliotis computes the diffusion coefficient of the exponential-correlation process as
$\int_0^\infty C = C(0)\,\tau_{\mathrm{cor}} = D/\alpha^2$, so for the OU kernel the two quantities
differ, $\Dstar = \int_0^\infty C = D/\alpha^2$, coinciding only at $\alpha=1$. (Eqn.~1.11 merely
sets $\alpha=D=1$ for a specific covariance computation, incidentally collapsing the distinction ---
it is not the locus of the conflation.) We keep $D$ for the noise strength and $\Dstar$ for the
transport coefficient throughout, so every rate constant carries its explicit $\alpha$-dependence.
This is what justifies estimating a covariance or a spectral density from a \emph{single} simulated
path and comparing it to a closed-form target --- the core verification methodology used throughout.

\subsection{Summary}

\begin{center}
\begin{tabular}{@{}p{3.1cm}p{4.3cm}p{4.3cm}@{}}
\toprule
& \textbf{Fourier / Bochner} & \textbf{Mercer / KL} \\
\midrule
Requires & weak stationarity & compact domain (continuous $\mathcal{C}$) \\
Object & spectral density $S(\omega)$ & eigenpairs $(\lambda_k,e_k)$ \\
Lives on & frequency axis $\mathbb{R}$ & discrete index $k\in\mathbb{N}$ \\
Exact coincidence & \multicolumn{2}{c}{on the torus only: $\lambda_k=\hat R(k)$} \\
\bottomrule
\end{tabular}
\end{center}

% =====================================================================
\section{Feature Ladder, Spine, and Cut Line}
% =====================================================================

\paragraph{Narrative spine vs.\ build DAG.} The \emph{narrative} spine --- the order in which the
units are best read --- is $0\to1\to2\to3\to4$, with Units~5, 6, and 7 hanging off it. That order is
pedagogical and is \emph{not} the dependency graph. The \emph{build DAG} --- what each unit actually
reuses, and therefore what each feature brief's ``depends on'' field records --- is
\[
1\Leftarrow 0,\quad 2\Leftarrow 0,\quad 3\Leftarrow 0,1,2,\quad 4\Leftarrow 0,1,\quad
5\Leftarrow 0,3,\quad 6\Leftarrow 1,2,4,\quad 7\Leftarrow 3.
\]
In particular \textbf{Unit~2 depends on Unit~0 only, not Unit~1} (the KL solver uses linear algebra,
not the spectral machinery); $1\leftrightarrow2$ is a narrative adjacency, not a build edge. Two of
these edges are \emph{preferential rather than strict}: Unit~4 samples OU via Unit~1's circulant
route \emph{preferred}, but Unit~0's Cholesky suffices (OU is well-conditioned); and Unit~7 borrows
the OU \emph{kernel} that first appears with Unit~1 and checks against Unit~3's exact stationary OU.
No edge points forward, so each unit is plannable against only the units below it.

\paragraph{Committed core and the cut line.} \textbf{Units 0--3 are shipped.} Units~4 and~5 are
committed core. \textbf{Unit~7 is the single cut line} --- a labelled forward pointer to Chapter~3
that nothing in Units 0--6 depends on, so it is the first to go if the debugging buffer is consumed;
\textbf{Unit~6 trims second}. (The source this plan replaces contradicted itself here, marking Unit~7
``the single cut line'' while also calling Unit~6 material to ``trim first alongside~7''; the
reconciliation is deliberate: 7 first because it is a pure forward pointer, 6 second because it is a
genuinely new kernel that the core does not require.) This protects depth over breadth: the budget
beyond the core is spent on \emph{more of Chapter~1 done thoroughly}, not on rushing ahead. See the
\hyperref[sec:budget]{Budget} for the hour split.

% =====================================================================
\section{Feature Briefs}
% =====================================================================

Each brief is orientation for the planner (\texttt{plan-and-dispatch}) that will decompose it into
commit plans, not a specification. Its fields: \textbf{Idea} (the one concept); \textbf{Proves} (the
claim and the exact analytic target it is checked against, with any external target flagged);
\textbf{Falsifier} (the hypothesis whose failure the unit is designed to exhibit --- the intent of
its negative control); \textbf{Depends on / enables} (its place in the build DAG); \textbf{Deliverable
shape} (the repository areas it touches and the class of artifact it emits, in prose --- never a call
signature or a tuned number); \textbf{Effort \& risk} (the estimate and the dominant overrun risk,
which points into the \hyperref[sec:risk]{risk register}); and \textbf{Status} on two axes ---
\emph{build} (shipped / planned) and \emph{cut} (core / trim-second / cut-line-trim-first). The exact
closed-form targets below are philosophy and are pinned here; the tolerances that gate them are
measurement and live in the code and the per-commit plans.

% ---------------------------------------------------------------
\subsection*{Unit 0 --- Covariance Core and Gaussian Sampling via Cholesky}
\addcontentsline{toc}{subsection}{Unit 0 --- Covariance Core (Cholesky)}
\begin{brief}
\item[Idea] A Gaussian process is fixed by its mean and covariance \pav{App.~B.5}; sampling is
applying a square root of $\Sigma$ (Cholesky, $\Sigma=LL^\top$, $X=Lz$).
\item[Proves] The empirical covariance of $N$ sample paths converges to the analytic
$\Sigma_{ij}=R(t_i,t_j)$ at the Monte-Carlo rate. The headline artifact is the fitted log--log slope
of $\|\widehat\Sigma_N-\Sigma\|_F$ against $N$, checked against the analytic $-\tfrac12$ and gated
stochastically (see \hyperref[sec:conv]{conventions}); the intercept is irrelevant, the exponent is
the object. A passing slope certifies the sampler and the estimator at once.
\item[Falsifier] The nugget is load-bearing. Factoring the \emph{same} singular Brownian-motion
$\Sigma$ used in the main check --- its $t=0$ row is identically zero, $R(0,s)=0$, so $\Sigma$ is
exactly singular --- with the nugget switched off throws; the throw disappears once the nugget is
restored. Reusing the main-check $\Sigma$ (rather than a separate pathological matrix) is the honest
control: the OU kernel is well-conditioned and would not throw.
\item[Depends on / enables] Root of the build DAG (depends on nothing). Enables Units~1--5: every
later unit draws paths through this.
\item[Deliverable shape] \texttt{gaussianprocess.jl} (the \texttt{GaussianProcess} type, grid
assembly of $\Sigma$, the empirical-covariance estimator) and \texttt{sampling.jl} (the jittered
Cholesky route). The \texttt{experiments/00\_covariance\_core} experiment emits a sample-paths figure
and the error-vs-$N$ log--log plot with the fitted slope annotated against the reference line, and
records the jitter and the RNG seed.
\item[Effort \& risk] $\approx 6$--$8$\,h (includes the repository skeleton). Dominant risk:
path-matrix orientation --- \texttt{empirical\_cov} expects paths as \emph{columns}
($n_{\text{grid}}\times N$); a transpose silently estimates the wrong matrix.
\item[Status] Build: \textbf{shipped}. Cut: \textbf{core}.
\end{brief}

% ---------------------------------------------------------------
\subsection*{Unit 1 --- Stationarity Gate: Spectral Density vs.\ Bochner (with Circulant Embedding)}
\addcontentsline{toc}{subsection}{Unit 1 --- Stationarity Gate (Spectral/Bochner)}
\begin{brief}
\item[Idea] Weak stationarity \pav{\S1.2} moves all second-order structure onto the frequency axis
(Bochner, \eqref{eq:bochner}); the Ornstein--Uhlenbeck / exponential kernel has a Lorentzian
spectrum,
\begin{equation}
R(\tau) = \frac{D}{\alpha}e^{-\alpha|\tau|} \quad\Longrightarrow\quad
S(\omega) = \frac{D}{\pi}\,\frac{1}{\omega^2+\alpha^2}.
\label{eq:lorentzian}
\end{equation}
\pav{Example 1.15} Circulant embedding runs the same Bochner/FFT machinery as an exact stationary
sampler.
\item[Proves] Two independent estimators of $S$ (an averaged Welch periodogram and the FFT of the
estimated autocorrelation) match the Lorentzian \eqref{eq:lorentzian} over the resolved band, and the
total-variance normalization holds: $\int_{\mathbb{R}}S\,d\omega = D/\alpha = R(0)$ under the
$1/2\pi$-on-$S$ convention of \eqref{eq:spectraldensity}. Circulant samples reproduce the analytic
covariance to floating-point scale. The headline artifact is the estimated-vs-analytic density on
log--log axes plus the pinned normalization scalar. \extnote{that the raw periodogram is a
statistically \emph{inconsistent} estimator of $S$ and must be smoothed or averaged --- a standard
time-series fact (Percival \& Walden; Brockwell \& Davis).}
\item[Falsifier] Two controls, both plotted. (i)~The raw periodogram overlaid on Welch: it stays
jagged and its variance does not shrink with record length at fixed $\omega$, while Welch converges
--- inconsistency made visible. (ii)~Dropping the $1/2\pi$ (engineering convention): the normalization
ratio lands on exactly $2\pi$ instead of $1$ --- the spurious factor made to appear on purpose. (This
README carries the $2\pi$ lesson for the whole repository.)
\item[Depends on / enables] $\Leftarrow$ Unit~0. Enables Units~3, 4, 6.
\item[Deliverable shape] \texttt{kernels.jl} (the exponential/OU kernel), \texttt{spectral.jl} (the
Welch and raw periodogram estimators, the Bochner forward/inverse transforms, and the total-power and
total-variance integrators), and \texttt{sampling.jl} (circulant embedding). The
\texttt{experiments/01\_spectral\_bochner} experiment emits the density plot with the resolved band
shaded, a raw-vs-Welch overlay, a circulant covariance-error panel, and the dropped-$1/2\pi$ control.
\item[Effort \& risk] $\approx 8$--$9$\,h --- the highest overrun risk in the project. Dominant risk:
the \hyperref[sec:risk]{FFT/normalization family} (one- vs two-sided spectra, angular vs ordinary
frequency, unnormalized FFTW, window power, circulant scaling) --- pin the normalization once so the
discrete $\int\widehat S\,d\omega\to R(0)$ and unit-test it.
\item[Status] Build: \textbf{shipped}. Cut: \textbf{core}.
\end{brief}

% ---------------------------------------------------------------
\subsection*{Unit 2 --- Karhunen--Lo\`eve via a Quadrature Eigenproblem (Torus Contrast + Truncation Cost)}
\addcontentsline{toc}{subsection}{Unit 2 --- Karhunen--Lo\`eve (Quadrature Eigenproblem)}
\begin{brief}
\item[Idea] The KL expansion \pav{\S1.5, Thm.~1.28} and Mercer's theorem give the operator's own
adapted-basis diagonalization; Brownian motion has a closed-form KL basis on $[0,1]$,
\begin{equation}
\lambda_n = \left(\frac{2}{(2n-1)\pi}\right)^{2}, \qquad
e_n(t) = \sqrt{2}\,\sin\!\left(\tfrac{1}{2}(2n-1)\pi t\right),
\label{eq:bm-kl}
\end{equation}
\pav{Example 1.29, eqn.~1.37} solved numerically by a Nystr\"om quadrature eigenproblem,
\emph{symmetrized} before the eigensolve.
\item[Proves] The numerical eigenpairs match the closed form \eqref{eq:bm-kl} and the eigenvalue decay
follows $\lambda_k\sim k^{-2}$; the trace identity $\sum_n\lambda_n = \int_0^T R(t,t)\,dt$ pins
assembly (this equals $T^2/2$ for Brownian motion --- \emph{not} $T\,R(0)=0$; the wrong target is
itself worth recording), with the stationary specialization $\sum_n\lambda_n=T\,R(0)$
\pav{Exercise 23} checked on a stationary kernel; and the torus coincidence $\lambda_k=\hat R(k)$
holds \emph{exactly} on the circle. The decay is a \emph{deterministic} slope (its residuals are
discretization misspecification, not sampling noise), gated only over the resolved range above the
discretization floor --- see \hyperref[sec:conv]{conventions}.
\item[Falsifier] Two, both latent in the design. (i)~Run the torus coincidence on $[0,T]$: it
\emph{fails}, because the interval eigenfunctions are Sturm--Liouville sines, not characters. (ii)~
Extend the eigenvalue plot \emph{below} the discretization floor: the $k^{-2}$ slope flattens --- the
resolution limit made visible rather than hidden.
\item[Depends on / enables] $\Leftarrow$ Unit~0 (not Unit~1). Enables Units~3, 6.
\item[Deliverable shape] \texttt{kl.jl} (quadrature nodes and weights, the symmetrized Nystr\"om
eigensolver, the diagonal-trace quadrature, the KL tail-energy), \texttt{sampling.jl} (the KL
truncation route), and \texttt{kernels.jl} (a torus kernel with a \emph{full} positive spectrum ---
derived in-house, since Pavliotis's Exercise~27 verifies $\lambda_k=\hat R(k)$ only for the
degenerate rank-2 kernel $\cos2\pi(t-s)$, insufficient for the torus contrast wanted here). The
\texttt{experiments/02\_kl\_quadrature} experiment emits the eigenfunctions, the eigenvalue decay with
the resolved window shaded and the floor annotated, side-by-side torus-Fourier vs.\ interval
Sturm--Liouville mode shapes, and the truncation-error-vs-$K$ and KL-vs-Cholesky cost curves.
\item[Effort \& risk] $\approx 10$--$11$\,h --- the deepest unit. Dominant risk: the Nystr\"om
symmetrization (the naive discretization is non-symmetric even though $R$ is) and the eigenvalue noise
floor --- both in the \hyperref[sec:risk]{risk register}.
\item[Status] Build: \textbf{shipped}. Cut: \textbf{core}.
\end{brief}

% ---------------------------------------------------------------
\subsection*{Unit 3 --- Reconciliation: The Process Zoo}
\addcontentsline{toc}{subsection}{Unit 3 --- Reconciliation (Process Zoo)}
\begin{brief}
\item[Idea] A small catalogue of processes --- Brownian motion $R(t,s)=\min(t,s)$ \pav{eqn.~1.21},
Brownian bridge $R(t,s)=\min(t,s)-ts$ \pav{eqn.~1.26 / Exercise 6}, stationary OU
$R(\tau)=\tfrac{D}{\alpha}e^{-\alpha|\tau|}$ \pav{Example 1.15} --- each just a choice of kernel. The
unit's one distinct idea is \emph{reconciliation}: the sampling routes and the two diagonalizations
are forced to agree on concrete processes, and \emph{distributional} (not merely moment) identities
are verified from finite samples. The sampling itself is reuse of Units~0/1/2.
\item[Proves] Three things. \emph{(a) Route equivalence:} Cholesky, KL, and circulant samples of OU
yield the same empirical covariance, judged against a \emph{split-half null} --- the same statistic
computed between two halves of a single route's samples, which is by construction a draw from the
``same law, different samples'' distribution and so calibrates the gate without a variance formula.
(This is essential because $\|\widehat\Sigma_A-\widehat\Sigma_B\|_F$ is a norm over many entries: it
concentrates about a strictly positive value, not zero, so a ``within $k$ SE of zero'' predicate gates
against the wrong null and passes routes that disagree by a factor of two.) \emph{(b) Distributional
identities}, in two steps that are ``the chain, not the conclusion'': the transformed process is
Gaussian \emph{by construction} (a deterministic linear time-change of a Gaussian process is Gaussian,
laws being closed under linear maps) --- a proof, not something the samples establish; App.~B.5 then
upgrades a match of the \emph{full} covariance into equality in law (e.g.\ $c^{-1/2}W(ct)\overset{d}{=}
W(t)$); and Cram\'er--Wold projections --- fixed linear functionals $a^\top X$, each KS-tested against
$N(0,a^\top\Sigma a)$ --- supply the evidence that would catch a non-Gaussian impostor carrying the
right covariance, which is precisely the evidence App.~B.5 alone cannot give. The KL coefficients are
checked the same way (empirical correlation of the $\xi_k$ is $\approx I$ and each $\xi_k$ is KS-tested
against $N(0,\lambda_k)$: uncorrelated \emph{and} jointly Gaussian $\Rightarrow$ independent, the KL
expansion's actual content). Bridge endpoints are pinned at $t=0,1$ via the covariance/Cholesky route,
because the time-change formula $B_t=(1-t)W(t/(1-t))$ cannot be evaluated at $t=1$. \emph{(c) The
cross-check:} on $[0,T]$ the OU KL eigenvalues converge, by the large-$T$ Toeplitz/Grenander--Szeg\H{o}
limit, to $\lambda_k\to\hat R(\omega_k)$, $\omega_k=k\pi/T$, where $\hat R(\omega)=2\pi\,S(\omega)$ is
the \emph{un-normalized} operator symbol --- the Nystr\"om eigenproblem carries no $1/2\pi$, so its
eigenvalues converge to $\hat R$, not to the Unit-1 density $S$ (which is $2\pi$ smaller). The
convergence is distributional, not uniform in $k$ (Szeg\H{o} controls linear statistics, weakest at
the spectral edges), so the gate is on the resolved-bulk gap shrinking with $T$ --- never a $\max_k$
over the whole spectrum --- and no theoretical decay exponent is claimed. This is a \emph{deterministic}
slope (the gap is computed against the analytic $\hat R$, not a Welch estimate). \extnote{the
Toeplitz/Szeg\H{o} asymptotic, the Cram\'er--Wold device, and the Kolmogorov--Smirnov statistic and its
finite-sample critical values are standard but external to Pavliotis.}
\item[Falsifier] Repeat the self-similarity identity with the \emph{wrong} exponent,
$c^{-1/3}W(ct)$: the empirical covariance no longer matches $\min(t,s)$ (it is off by $c^{1/3}$),
which is what makes the correct $c^{-1/2}$ meaningful. Fully self-contained.
\item[Depends on / enables] $\Leftarrow$ Units~0, 1, 2 (the three sampling routes). Enables Units~5
and~7.
\item[Deliverable shape] Reuses \texttt{sample\_cholesky} / \texttt{sample\_kl} /
\texttt{sample\_circulant\_embedding}, \texttt{empirical\_cov} (also per half, for the split-half
null), \texttt{nystrom\_eigen} (the cross-check eigenvalues), the three catalogue kernels including a
new Brownian-bridge kernel, and Welch (only for the visual $S$-overlay in the artifact, never the
gate). It introduces \emph{one} new \texttt{src/} module, \texttt{gof.jl} (a Kolmogorov--Smirnov
statistic): deliberately \emph{not} an operation on $\mathcal{C}$ --- the one considered exception to
the organizing principle --- placed in \texttt{src/} rather than the experiment folder because two
units need it (Cram\'er--Wold here, Donsker rates in Unit~5) and \texttt{src/} is the tier that gets a
deterministic test. The \texttt{experiments/03\_process\_zoo} experiment emits a three-panel portrait
per process, a route-equivalence panel plotting the pairwise distances against the split-half null
band, a Cram\'er--Wold per-projection panel, the KL-coefficient correlation/normality check, the
pinned bridge endpoints, the cross-check $\log g(T)$-vs-$\log T$ plot with its fitted slope, and the
wrong-exponent control.
\item[Effort \& risk] $\approx 8$\,h --- almost entirely reuse; the added hour is the Cram\'er--Wold
projections, the split-half calibration, and \texttt{gof.jl} with its deterministic test. Dominant
risk: the \hyperref[sec:risk]{estimator-noise-floor family} applied to the cross-check bulk gap ---
gate on the bulk, never on a $\max_k$ where the edge eigenvalues Szeg\H{o} does not control can stall
the decay.
\item[Status] Build: \textbf{shipped}. Cut: \textbf{core}.
\end{brief}

% ---------------------------------------------------------------
\subsection*{Unit 4 --- Ergodicity as Loop-Closer}
\addcontentsline{toc}{subsection}{Unit 4 --- Ergodicity as Loop-Closer}
\begin{brief}
\item[Idea] The ergodic properties of stationary processes \pav{\S1.2}: the $L^2$ law of large
numbers \pav{Prop.~1.16, \eqref{eq:ergodic}}, valid under $C\in L^1(0,\infty)$; the finite-$T$
variance identity \pav{Lemma~1.17, \eqref{eq:ergodic-rate}}; and the Green--Kubo formula
\eqref{eq:greenkubo}.
\item[Proves] For stationary OU, the variance of the time-average obeys
$\operatorname{Var}\!\big(\tfrac1T\int_0^T X_s\,ds\big)\to \tfrac{2\Dstar}{T}=\tfrac{2D}{T\alpha^2}$
--- both the slope $-1$ and the constant, the latter carrying $\alpha$ --- and the integrated process
grows as $\mathbb{E}\big(\int_0^t X_s\,ds\big)^2\to 2\Dstar t = \tfrac{2D}{\alpha^2}t$. The headline
artifact is the variance-vs-$T$ log--log slope gated stochastically against $-1$, together with the
fitted constant against $2\Dstar$ (the constant is what pins $\alpha$). \emph{This is the
methodological keystone}: it is the justification for the whole project's one-path-vs-analytic
methodology used in Units~0--3.
\item[Falsifier] Feed the \emph{same} time-average-variance estimator a synthetic stationary process
with a \emph{non-integrable} correlation, $C(u)\propto(1+u)^{-1/2}$ (so $C\notin L^1$): the variance
decays slower than $1/T$ and the fitted slope departs from $-1$ --- Prop.~1.16's $L^1$ hypothesis
shown load-bearing. Self-contained, with no forward reference to fBm (Unit~6 runs the genuine-process
version).
\item[Depends on / enables] $\Leftarrow$ Units~0 and~1 --- reuses Unit-1 circulant (preferred) or
Unit-0 Cholesky for OU paths. Enables Unit~6.
\item[Deliverable shape] A new \texttt{ergodic.jl} (a running time-average along one path, the
time-average variance across an \emph{ensemble} of length-$T$ paths, the integrated mean-square
displacement, and the Green--Kubo coefficient $\Dstar=D/\alpha^2$), plus the reused OU sampler. The
\texttt{experiments/04\_ergodicity} experiment emits the variance-vs-$T$ log--log with the fitted
$1/T$ slope and constant, the MSD-vs-$t$ with fitted slope, the running-average path with a
$\pm\sqrt{2\Dstar/T}$ band, and the non-integrable-$C$ control.
\item[Effort \& risk] $\approx 6$\,h. Dominant risk: confusing the time-average with its variance ---
one long path gives a single realization, not the variance, which needs an ensemble (a path
\emph{matrix}), so compute the $\pm\sqrt{2\Dstar/T}$ band with $\Dstar=D/\alpha^2$, not $D$.
\item[Status] Build: \textbf{planned}. Cut: \textbf{core}.
\end{brief}

% ---------------------------------------------------------------
\subsection*{Unit 5 --- Brownian Motion as a Scaling Limit (Random Walk $\to$ BM)}
\addcontentsline{toc}{subsection}{Unit 5 --- Brownian Motion as a Scaling Limit}
\begin{brief}
\item[Idea] Brownian motion as the weak (Donsker) limit of a rescaled random walk \pav{\S1.3}: for
iid mean-zero, variance-one increments the interpolated, rescaled walk converges weakly to a standard
Brownian motion. This is BM's \emph{construction} --- a statement about laws on path space,
complementary to the covariance-operator picture and deliberately off the $\mathcal{C}$-spine.
\item[Proves] The finite-dimensional distributions and a path functional (the running maximum)
converge to the Gaussian/BM target: the one-time marginal $\Rightarrow N(0,1)$, the two-time
covariance $\to\min(s,t)$. The \emph{marginal-convergence rate is law-dependent}: Berry--Esseen bounds
the Kolmogorov--Smirnov statistic by $O(n^{-1/2})$, but that rate is achieved only when the skewness is
nonzero --- for a \emph{smooth} symmetric increment the leading $n^{-1/2}$ Edgeworth term vanishes,
leaving the $O(n^{-1})$ kurtosis term. A \emph{lattice} symmetric law, however, carries a separate,
skewness-independent Esseen lattice-discreteness term that is itself $O(n^{-1/2})$ and does not vanish.
So the target is $\propto n^{-1/2}$ for the \emph{exponential} (skewed) increment \emph{and} for
\emph{Rademacher} (symmetric but lattice, via discreteness), and $\propto n^{-1}$ only for the
smooth-symmetric \emph{uniform} increment --- three mechanisms, two rates. The headline artifact is
the KS-vs-$n$ log--log plot with one line per law, each slope gated stochastically. \extnote{Donsker's
invariance principle / the functional CLT; Berry--Esseen for the $n^{-1/2}$ bound and the Edgeworth
expansion for the $n^{-1}$ symmetric-law rate. Standard, external to the second-order framing;
high confidence.}
\item[Falsifier] Add an \emph{infinite-variance} increment (symmetric Pareto, tail index $<2$): the
rescaled walk does not converge to $N(0,1)$ at all --- the marginal departs from the QQ line, its
empirical variance \emph{grows} with $n$ instead of settling at $1$, and $n^{-1/2}$ is simply the
wrong normalization --- showing Donsker's finite-variance hypothesis is load-bearing.
\item[Depends on / enables] $\Leftarrow$ Units~0 (the empirical covariance) and~3 (\texttt{gof.jl}'s
KS statistic). Enables nothing.
\item[Deliverable shape] \emph{No new \texttt{src/} code} --- the rescaled-walk builders live in the
experiment, weak convergence of laws not being an operation on $\mathcal{C}$; the one library call is
\texttt{gof.jl}'s \texttt{ks\_statistic}, introduced by Unit~3 and reused rather than duplicated. The
\texttt{experiments/05\_bm\_scaling\_limit} experiment emits overlaid rescaled paths at increasing
$n$, a marginal histogram/QQ against $N(0,1)$, the KS-vs-$n$ log--log with one line per increment law,
the empirical-vs-$\min(s,t)$ covariance, the running-maximum across laws, and the infinite-variance
control.
\item[Effort \& risk] $\approx 4$--$5$\,h (plus a small buffer for the law-dependent-rate surprise).
Dominant risk: expecting $n^{-1/2}$ for the symmetric laws --- a ``too steep'' $-1$ slope is correct,
not a bug.
\item[Status] Build: \textbf{planned}. Cut: \textbf{core}.
\end{brief}

% ---------------------------------------------------------------
\subsection*{Unit 6 --- Fractional Brownian Motion}
\addcontentsline{toc}{subsection}{Unit 6 --- Fractional Brownian Motion}
\begin{brief}
\item[Idea] A one-exponent generalization of Brownian motion tuning both self-similarity and increment
correlation \pav{Def.~1.26, Prop.~1.27}, with covariance
\begin{equation}
R_H(t,s) = \tfrac{1}{2}\left(|t|^{2H} + |s|^{2H} - |t-s|^{2H}\right).
\label{eq:fbm}
\end{equation}
\pav{eqn.~1.27}
\item[Proves] The empirical covariance matches $R_H$ for several $H$; the self-similarity identity
$W^H_{ct}\overset{d}{=}c^{H}W^H_t$ \pav{eqn.~1.28} holds as a covariance identity; and the eigenvalue
decay of $\mathcal{C}_H$ contrasts with Brownian motion's $k^{-2}$. The headline artifacts are the
covariance and self-similarity matches together with the eigenvalue-decay overlay --- and, as the
control below, the ergodic breakage at $H>\tfrac12$.
\item[Falsifier] The full fBm version of Unit~4's control: apply the time-average-variance estimator
at $H=0.8$. The increments have non-summable correlations ($C\notin L^1$), so the $1/T$ ergodic rate
breaks --- the first genuine process where Unit~4's machinery fails, run rather than asserted.
\item[Depends on / enables] $\Leftarrow$ Units~1 (circulant sampler, with the fBm
minimal-nonnegative-embedding caveat), 2 (\texttt{nystrom\_eigen} for the decay contrast), and 4
(the time-average-variance estimator for the breakage control); also Unit~0's Cholesky. Enables
nothing.
\item[Deliverable shape] \texttt{kernels.jl} gains the fBm kernel; the unit reuses the circulant and
Cholesky samplers, \texttt{nystrom\_eigen}, and \texttt{time\_average\_variance}. The
\texttt{experiments/06\_fbm} experiment emits paths at $H\in\{0.3,0.5,0.8\}$ (rougher to smoother),
covariance heat-maps, the self-similarity scaling check, the $\mathcal{C}_H$ eigenvalue-decay overlay
against BM, and the $H=0.8$ ergodic-breakage panel.
\item[Effort \& risk] $\approx 5$--$6$\,h. Dominant risk: circulant non-negativity fails for fBm ---
the plain symmetric embedding can have negative eigenvalues, so enlarge to the minimal-nonnegative
(Dietrich--Newsam) embedding \extnote{fBm circulant embedding (Dietrich--Newsam) and the
nonnegative-embedding caveat are standard but external to Pavliotis.} --- part of the
\hyperref[sec:risk]{FFT/normalization and noise-floor families}.
\item[Status] Build: \textbf{planned}. Cut: \textbf{cuttable --- trim second}.
\end{brief}

% ---------------------------------------------------------------
\subsection*{Unit 7 --- SDE Bridge: OU Invariant Law and Strong vs.\ Weak Convergence}
\addcontentsline{toc}{subsection}{Unit 7 --- SDE Bridge (Strong vs.\ Weak Convergence)}
\begin{brief}
\item[Idea] The dynamical (SDE) view that bridges to Chapter~3. The OU SDE with \emph{additive} noise,
\begin{equation}
dX_t = -\alpha X_t\,dt + \sqrt{2D}\,dW_t,
\label{eq:ou-sde}
\end{equation}
has stationary law $N(0,D/\alpha)$ with covariance $\tfrac{D}{\alpha}e^{-\alpha|\tau|}$ \pav{cf.\
Ch.~2, Example 2.5} --- the $\sqrt{2D}$ makes the stationary variance agree with the Unit-3 kernel ---
and its Euler--Maruyama discretization. The strong/weak-convergence \emph{distinction} is demonstrated
on geometric Brownian motion with \emph{multiplicative} noise,
\begin{equation}
dY_t = \mu Y_t\,dt + \sigma Y_t\,dW_t, \qquad
Y_t = Y_0\exp\!\big((\mu-\tfrac12\sigma^2)t + \sigma W_t\big),
\label{eq:gbm}
\end{equation}
which has a closed-form solution against which both errors can be measured.
\item[Proves] The EM OU invariant law converges to $N(0,D/\alpha)$ with covariance $\to
\tfrac{D}{\alpha}e^{-\alpha|\tau|}$ (checked against Unit~3's exact stationary OU); and for GBM the
strong error $\mathbb{E}|Y^{\Delta t}_T-Y_T|\propto\Delta t^{1/2}$ (slope $\tfrac12$) while the weak
error $|\mathbb{E}f(Y_T^{\Delta t})-\mathbb{E}f(Y_T)|\propto\Delta t$ (slope $1$). The headline
artifact is the log--log plot of both GBM errors with their slopes gated stochastically \emph{and}
shown visibly distinct --- the $0.5$/$1.0$ gap being a multiplicative-noise phenomenon invisible to
deterministic Taylor analysis. \extnote{the OU/GBM SDEs and the Euler--Maruyama/Milstein strong and
weak convergence orders are Chapter-3 / standard-SDE-numerics material (Kloeden \& Platen); used here
only as an explicitly labelled forward pointer.}
\item[Falsifier] The most striking falsifier in the project: recompute the GBM strong error with
\emph{independent} coarse and fine Brownian paths instead of the shared one --- the ``strong error''
stops decaying and plateaus at an $O(1)$ constant, showing strong convergence is a statement about a
\emph{shared} noise realization. A near-free contrast reinforces it: the same integrator on OU
(additive noise) shows strong order $1.0$ --- no gap --- because the Milstein correction vanishes for
additive noise, which is precisely \emph{why} the strong/weak gap must be exhibited on multiplicative
noise; on GBM the same integrator shows $0.5$.
\item[Depends on / enables] $\Leftarrow$ Unit~3 (the exact stationary OU the invariant-law check is
compared against). Enables nothing --- it is the forward pointer.
\item[Deliverable shape] A new \texttt{sde.jl} (an Euler--Maruyama step, an additive-noise OU step, a
GBM closed-form reference, a path simulator from a step rule, a driving Brownian path with its
increments, and a coarsen-increments helper that shares one driving path across resolutions). The
\texttt{experiments/07\_sde\_bridge} experiment emits the OU invariant-law check (histogram/QQ and
covariance against the exact kernel), the GBM strong (shared path) and weak (many paths) errors
log--log with both slopes, the un-coupled strong-error plateau, and the OU-vs-GBM strong-order
contrast.
\item[Effort \& risk] $\approx 8$\,h. Dominant risk: the coupling trap (strong-error checks require
the \emph{same} driving path at coarse and fine steps --- generate the fine path and subsample it; the
un-coupled version is the control, not the method) and the weak-error Monte-Carlo floor --- both in
the \hyperref[sec:risk]{risk register}.
\item[Status] Build: \textbf{planned}. Cut: \textbf{cut line --- trim first} (a pure Chapter-3
forward pointer; nothing in Units 0--6 depends on it).
\end{brief}

% =====================================================================
\section{Repository Architecture and Rationale}
% =====================================================================

The repository separates a small, tested \textbf{library} (\texttt{src/}, organized by
\emph{operation on $\mathcal{C}$}) from a narrative, reproducible \textbf{experiment gallery}
(\texttt{experiments/}, organized by \emph{unit}). This separation is the central design decision.
Units 00--03 are shipped; 04--07 are planned and are \emph{not} pre-stubbed.

\begin{verbatim}
StochasticProcesses.jl/
|-- Project.toml, Manifest.toml     # both committed: exact reproducibility
|-- README.md                       # through-line + gallery index
|-- src/
|   |-- StochasticProcesses.jl      # module definition, flat re-exports
|   |-- kernels.jl                  # R(t,s): min, min-ts, exp/OU, periodic; (+fBm, Unit 6)
|   |-- gaussianprocess.jl          # GP type, grid assembly, empirical covariance
|   |-- sampling.jl                 # the square roots: Cholesky / KL / circulant (FFTW)
|   |-- spectral.jl                 # Welch & raw periodogram, Bochner transforms, integrators
|   |-- kl.jl                       # Nystrom quadrature eigenproblem (symmetrized)
|   |-- gof.jl                      # KS statistic; the considered exception -- see below
|   |-- ergodic.jl                  # time-average estimators, Green-Kubo   (Unit 4, planned)
|   `-- sde.jl                      # Euler-Maruyama, coupling helpers      (Unit 7, planned)
|-- test/                           # analytic identities, tight tolerances (CI runs this)
|   `-- runtests.jl
|-- experiments/                    # the gallery -- one folder per unit
|   |-- 00_covariance_core/     {run.jl, README.md, figures/}   # shipped
|   |-- 01_spectral_bochner/    {run.jl, README.md, figures/}   # shipped
|   |-- 02_kl_quadrature/       {run.jl, README.md, figures/}   # shipped
|   |-- 03_process_zoo/         {run.jl, README.md, figures/}   # shipped
|   |-- 04_ergodicity/          ...                             # planned
|   |-- 05_bm_scaling_limit/    ...                             # planned
|   |-- 06_fbm/                 ...                             # planned
|   `-- 07_sde_bridge/          ...                             # planned (Ch.3 pointer)
|-- experiments/Project.toml, Manifest.toml   # separate env that `dev`s the package
`-- .github/workflows/ci.yml
\end{verbatim}

\textbf{Rationale.} \texttt{src/} is organized by \emph{operation on $\mathcal{C}$} --- kernels, the
two diagonalizations, the square roots, the ergodic loop --- rather than by named process (Brownian
motion, OU, \ldots). A process is nothing more than a choice of kernel; organizing by operation makes
the module structure a faithful expression of the through-line \emph{for the second-order theory of
Chapter~1}. Two units sit deliberately at the edge of that spine and are placed among
\texttt{experiments/} accordingly: Unit~5 (random walk $\to$ BM) is weak convergence of \emph{laws},
not an operation on $\mathcal{C}$, and its builders live entirely in its experiment folder; Unit~7
(SDE bridge) is a dynamical, generator-flavored Chapter-3 pointer whose only library home is
\texttt{sde.jl}.

One module, \texttt{gof.jl}, is a \emph{considered exception} to the organizing principle: a
Kolmogorov--Smirnov statistic is not an operation on $\mathcal{C}$ by any reading, but Units~3 and~5
both need one --- for Cram\'er--Wold projections and for Donsker rates respectively --- and the
alternative is the same few lines duplicated across two experiment folders where the testing
convention gives them no deterministic test at all. It is placed in \texttt{src/} for the test
coverage and the single definition, and the principle is stated as \emph{bending} here rather than
quietly redefined to accommodate it.

\textbf{The rationale's limit, stated with the rationale.} The covariance-operator spine is the right
organizing principle precisely because Chapter~1 is the Gaussian/second-order chapter. Adding
later-chapter material never forces a rewrite of existing files --- one adds modules to \texttt{src/}
and numbered folders to \texttt{experiments/} --- but the covariance operator ceases to be the hub.
Chapter~2 (diffusions, Markov generators, the Kolmogorov equations) and Chapter~4 (Fokker--Planck) are
centered on the \emph{generator}/transition semigroup on generally non-Gaussian processes; they
introduce a second, generator-centered spine \emph{alongside} $\mathcal{C}$, rather than more
operations on it. So the code accretes without restructuring, but the single-hub \emph{narrative} is
Chapter-1-specific and should not be over-sold as the framework's permanent organizing idea. Each
experiment folder remains self-contained and documents concept~$\to$~implementation~$\to$~theoretical
check~$\to$~result, which directly serves the goal of leaving durable evidence of understanding.

% =====================================================================
\section{Cross-Cutting Conventions}
\label{sec:conv}
% =====================================================================

These are load-bearing for reproducibility; violating one silently corrupts every downstream number.
The concrete values (seeds, jitter, grid sizes, tolerances) are recorded in each experiment, not here.

\begin{itemize}[leftmargin=1.4cm]
  \item \textbf{Seeding.} Use \texttt{StableRNGs.jl}, never Julia's default global RNG. \extnote{the
  default stream is not guaranteed stable across Julia versions, so seeded Monte Carlo checks can
  silently drift; a Julia-ecosystem fact, not a Pavliotis one.} Pass an explicit
  \texttt{StableRNG(seed)} to every stochastic routine and record the seed. Do not reorder RNG draws
  within an experiment --- demo paths and the ladder pull from one stream, so reordering changes every
  committed number.
  \item \textbf{Numerical positive-definiteness (Cholesky).} Form $\Sigma$ on the grid and factor
  $\Sigma+\varepsilon I$ with a small jitter (nugget) $\varepsilon$, and \emph{report} it: too small
  gives an indefinite throw, too large biases the covariance. Prefer the KL or circulant route for the
  smoothest kernels where even jittered Cholesky loses accuracy.
  \item \textbf{Spectral normalization.} Pin the periodogram and Bochner-FFT normalizations so the
  discrete $\int\hat S\,d\omega\to R(0)$, matching the $1/2\pi$-on-$S$ convention of
  \eqref{eq:spectraldensity}; a mismatched normalization is the most likely source of a spurious
  factor of $2\pi$. Document the chosen normalization (and the circulant-embedding scaling) in the
  relevant experiment README.
  \item \textbf{Two verification tiers.}
  \begin{enumerate}
    \item \texttt{test/}: deterministic analytic identities --- kernel symmetry/PSD, the trace
    identity in its general form $\sum_n\lambda_n=\int_0^T R(t,t)\,dt$ (plus the stationary
    specialization $T\,R(0)$), Brownian-motion eigenvalues vs.\ the closed form \eqref{eq:bm-kl} ---
    with tight absolute tolerance, over the resolved eigenvalue range only. This is what CI runs.
    \item \texttt{experiments/}: Monte Carlo checks with tolerance a small multiple of the estimated
    standard error, seeded so results are deterministic given the seed.
  \end{enumerate}
  \item \textbf{Tests must bite.} Prefer hand-computed targets over re-running the function's own
  formula, and non-square / asymmetric fixtures so shape, orientation, and one-sided-implementation
  bugs cannot hide (e.g.\ the KS statistic is exercised on skewed fixtures, since a symmetric case does
  not discriminate a one-sided implementation).
  \item \textbf{Rate checks as the primary evidence.} A measured \emph{slope} on a log--log plot is
  the headline artifact for every unit, and is stronger evidence of correct implementation than a
  single-point match. The two classes of slope are gated differently, and \emph{the distinction is not
  whether the process is random but whether the regression residuals are sampling noise}:
  \begin{itemize}[leftmargin=1.0cm]
    \item \emph{Stochastic slopes} (the Monte-Carlo covariance error $\propto N^{-1/2}$; the
    Euler--Maruyama strong/weak orders on GBM; the ergodic variance $\propto 1/T$; the Donsker
    KS-vs-$n$ rates of Unit~5): fit on log--log axes and gate against the theoretical value to within a
    small multiple of the fitted slope's own standard error --- the estimator variance differs sharply
    between checks, so a single fixed absolute tolerance would be simultaneously too loose in one place
    and too tight in another.
    \item \emph{Deterministic slopes} (the KL eigenvalue-decay slope of Unit~2; the Toeplitz bulk-gap
    slope of Unit~3): the regression residuals are model misspecification --- the asymptotic $k^{-2}$
    plus the discretization floor for Unit~2, the Szeg\H{o} asymptotic plus quadrature error for
    Unit~3 --- \emph{not} sampling noise, so a regression standard error is not a sampling standard
    error and an SE-based gate is not meaningful. Fit only within the resolved window and gate with a
    fixed absolute margin: for Unit~2 against the theoretical exponent $-2$; for Unit~3, where no
    defensible theoretical exponent exists, against a fixed negative threshold, reporting the fitted
    value rather than claiming it. Unit~3's gap is deterministic precisely because it is computed
    against the analytic symbol rather than a Welch estimate.
  \end{itemize}
  \item \textbf{Continuous integration.} Run \texttt{test/runtests.jl} on push via GitHub Actions ---
  fast, guarding the deterministic analytic core. Do \emph{not} run the slower Monte Carlo experiments
  in CI; run those locally and commit the resulting figures directly. Experiments are headless
  (\texttt{ENV["GKSwstype"]="100"} before \texttt{gr()}).
  \item \textbf{Environment pinning.} Commit \texttt{Manifest.toml} alongside \texttt{Project.toml}
  (for \emph{both} the root package env and the separate \texttt{experiments/} env, which \texttt{dev}s
  the package via a relative path) so \texttt{Pkg.instantiate()} reproduces exact versions. Manifest
  pinning plus explicit RNG seeding gives full end-to-end reproducibility of every figure.
\end{itemize}

% =====================================================================
\section{Consolidated Risk Register}
\label{sec:risk}
% =====================================================================

Three failure families recur across units; internalize them once rather than per unit.

\begin{itemize}[leftmargin=1.4cm]
  \item \emph{The FFT/normalization family} (Units 1, 6, 7): one- vs two-sided spectra; angular
  ($\omega$) vs ordinary ($f=\omega/2\pi$) frequency, worth a factor $2\pi$; FFTW returns unnormalized
  transforms (restore $\Delta t$, $1/N$); Welch divides by the window power $\sum_n w_n^2$, not $N$;
  the circulant-embedding scaling is validated once, in Unit~1. Single guard: pin every normalization
  so the discrete $\int\widehat S\,d\omega\to R(0)$ and unit-test it.
  \item \emph{The estimator noise floor} (Units 2, 7, and the raw periodogram of Unit~1): every rate
  check has a level below which the estimator's own variance or discretization dominates and the slope
  flattens --- the Nystr\"om eigenvalue floor, the weak-error Monte-Carlo floor, and the inconsistent
  raw periodogram are the \emph{same} phenomenon. Single guard: identify the floor, fit slopes only
  above it, and state the resolved range.
  \item \emph{The strong-convergence coupling} (Unit 7): strong-error checks require the \emph{same}
  driving Brownian path at coarse and fine steps --- generate the fine path first and subsample it.
  The un-coupled version is a negative control, not the method.
\end{itemize}

% =====================================================================
\section{Budget}
\label{sec:budget}
% =====================================================================

At the working budget of $15$\,h/week ($6\times2.5$\,h) against a $6$--$8$ week horizon, i.e.\
$90$--$120$\,h, the summed per-unit estimates total $\approx 55$--$61$ core hours; the remainder is
deliberate headroom for per-unit spec-writing and review of all generated code, plus a debugging
buffer concentrated on the units most likely to overrun for someone newer to stochastic numerics:
\textbf{Units 1, 2, 6, and 7}. The risk in Units~1, 6, and~7 is spectral/FFT normalization; the risk
in Unit~2 is the Nystr\"om symmetrization and the eigenvalue noise floor. One further hidden time-sink:
the Unit-5 marginal-convergence rate is law-dependent, so a fitted slope that is not $n^{-1/2}$ can be
correct rather than a bug.

Reconciled against shipped reality: \textbf{Units 0--3 are landed}, accounting for roughly $32$--$36$
of the core hours; the \textbf{remaining} committed work (Units 4--7) is roughly $23$--$25$ hours, so
the split reconciles with the total rather than merely relabeling it. Everything through Unit~6 is
firm; the SDE bridge (Unit~7) is the single cut line, with Unit~6 trimmed second, if the buffer is ever
consumed --- nothing in Units 0--6 depends on Unit~7. This spends the extra budget over a leaner plan
on \emph{more of Chapter~1 done thoroughly} (the scaling-limit construction, fBm, the truncation-cost
study, circulant embedding), consistent with the stated priority of depth over breadth.

% =====================================================================
\section*{Sources}
\addcontentsline{toc}{section}{Sources}
% =====================================================================

Call signatures, implementations, and numerical tolerances are \emph{not} recorded in this plan: the
shipped code under \texttt{src/} is the source of truth for the former, and each feature brief is
handed to \texttt{plan-and-dispatch}, whose per-commit plans own the contracts and the gate values.
(The previous revision of this document carried an appendix of illustrative code stubs and per-unit
signature tables; those were deleted as downstream content that competes with the real source of
truth.)

\begin{itemize}[leftmargin=1.4cm]
  \item G.~A.~Pavliotis, \emph{Stochastic Processes and Applications}. The project's copy is the
  \emph{November 11, 2015 draft}; all section, theorem, and equation numbers cited here (Chapter~1,
  plus \S B.5 and a single Chapter-2 pointer, Example~2.5) refer to that draft and were checked
  directly against it. Its numbering differs from the published Springer edition. In particular,
  \textbf{Mercer's theorem is stated inline and unnumbered in \S1.5} (there is no ``Theorem~1.18'';
  ``Example~1.18'' is the exponential-covariance / diffusion-coefficient example, where
  $\int_0^\infty C = D/\alpha^2$ is computed explicitly), and the ergodic result appears as
  \emph{Proposition~1.16} (its proof header inconsistently reads ``Theorem~1.16'').
  \item External references, flagged inline where used and not verified against Pavliotis: Grenander \&
  Szeg\H{o}, \emph{Toeplitz Forms and Their Applications} (asymptotic KL/spectral-density coincidence
  off the torus); D.~B.~Percival \& A.~T.~Walden, \emph{Spectral Analysis for Physical Applications},
  or P.~J.~Brockwell \& R.~A.~Davis, \emph{Time Series: Theory and Methods} (periodogram
  inconsistency); C.~R.~Dietrich \& G.~N.~Newsam, and A.~T.~A.~Wood \& G.~Chan (circulant embedding of
  stationary Gaussian fields); P.~E.~Kloeden \& E.~Platen, \emph{Numerical Solution of Stochastic
  Differential Equations} (Euler--Maruyama/Milstein strong and weak orders, including strong order~1
  for additive noise); W.~Feller, \emph{An Introduction to Probability Theory and Its Applications,
  Vol.~II} (the Edgeworth expansion / skewness-dependent Berry--Esseen rate of Unit~5); the functional
  central limit theorem / Donsker's invariance principle (Unit~5); and P.~Billingsley,
  \emph{Probability and Measure} (the Cram\'er--Wold device) together with any standard
  nonparametric-statistics reference for the Kolmogorov--Smirnov statistic and its finite-sample
  critical values (both used for the Unit-3 distributional identities).
\end{itemize}

\end{document}
